\section{Information Protection and Related Work}\label{sec:related}
\subsection{Inference through operational disclosures}
The inference and aggregation problem arises when combinations of accessible data reveal protected information~\cite{lunt1989aggregation,farkas2002inference}. In power infrastructure, public prices have been used to identify non-public network structure~\cite{kekatos2014topology}. Smart-meter privacy research likewise shows that operational measurements contain information beyond their immediate measurement purpose~\cite{mcdaniel2009smartgrid,giaconi2018smartmeter}. These findings motivate examining confidentiality at the level of the information carried by a release.

The protection objective here is field-level disclosure review. The auditor specifies operational targets and the fields available before those targets are released, then assesses their combinations. Record-level mechanisms such as differential privacy and $k$-anonymity~\cite{dwork2014dp,sweeney2002kanonymity} address related privacy objectives under release transformations; the present task is to measure the inferability associated with candidate operational fields.

\subsection{From relationships to field grades}
Inference-graph assessment~\cite{hu2026ignn} represents fields as nodes, encodes physical or statistical inference relations as edges, and propagates sensitivity from intrinsic scores. This approach incorporates domain knowledge and exposes chains of dependence. Our audit evaluates field sets directly, allowing a combination to be informative even when its members have weak individual links to the target. The score is defined by the inference outcome, and the background combinations provide evidence for its interpretation.

Feature-importance methods also derive field scores from set values. Shapley-based measures and SAGE average marginal contributions over contexts~\cite{shapley1953value,lundberg2017shap,covert2020sage,williamson2020spvim}; LOCO and permutation measures emphasize the context of a fitted model~\cite{lei2018loco,breiman2001rf,fisher2019mcr,hooker2021unrestricted}. The marginal contribution index takes the maximum contextual gain~\cite{catav2021mci}, consistent with an audit concerned about an informative admissible background. We specialize this quantity to a public baseline and a background budget, and report the amplification and witnesses needed for release review.

\subsection{Making repeated assessment feasible}
Evaluating many field sets creates the same refitting cost encountered in variable-importance estimation. Related approaches include linearized and warm-started models~\cite{gao2022lazyvi,sun2025warmstart}, masked surrogates~\cite{covert2021explaining,jethani2022fastshap}, and in-context prediction with tabular foundation models~\cite{hollmann2025tabpfn}. Masked reconstruction~\cite{yoon2020vime,he2022mae} and closed-form adaptation on shared features~\cite{bertinetto2019r2d2,lee2019metaoptnet} provide building blocks for our estimator. Their role is to make the infrastructure audit affordable to repeat; selected risk evidence is verified by dedicated retraining.
